\documentclass[11pt]{article}
\usepackage[margin=1in]{geometry}
\usepackage{booktabs}
\usepackage{longtable}
\usepackage{hyperref}
\usepackage{xcolor}
\usepackage{tabularx}
\usepackage{enumitem}
\usepackage{parskip}

\hypersetup{colorlinks=true,linkcolor=blue,urlcolor=blue,citecolor=blue}

\title{Replication Report: Chan et al.\ (2020)\\
\large ``AbGRI4, a novel antibiotic resistance island in multiply antibiotic-resistant \textit{Acinetobacter baumannii} clinical isolates''}
\author{Ollie (OpenClaw AI subagent) \\ BVBRC Replication Project (target BVBRC-69)}
\date{2026-07-03}

\begin{document}
\maketitle

\begin{center}
\fbox{\parbox{0.85\textwidth}{\centering
\textbf{Verdict:} \textcolor{green!50!black}{\textbf{REPLICATED}}\\
All 11 testable structural and molecular claims were independently reproduced on the paper's own public NCBI deposits (CP035043--CP035053) plus AB0057 and ATCC~17978 as external comparators, using only free/open tools on \texttt{uicgpu}.
}}
\end{center}

\section{Paper}
Chan AP, Choi Y, Clarke TH, Brinkac LM, White RC, Jacobs MR, Bonomo RA, Adams MD, Fouts DE.
\textit{Journal of Antimicrobial Chemotherapy} 75(10):2760--2768, 2020.
DOI: \href{https://doi.org/10.1093/jac/dkaa266}{10.1093/jac/dkaa266}. PMID:~32681170. PMCID:~PMC7556812. Open access (\copyright~OUP / OA via PMC).

Chan et al.\ sequenced four multi-drug-resistant \textit{A.\ baumannii} isolates (ABUH763, ABUH773, ABUH793, ABUH796) with combined Illumina + Oxford Nanopore reads and closed chromosomes and plasmids (CP035043--CP035053). Three of four (ABUH763, ABUH793, ABUH796; Pasteur ST2 / Oxford ST281, clade F) carry a novel $\sim$6.8-kb, IS26-bounded class 1 integron resistance island --- named \textbf{AbGRI4} --- inserted at a previously undescribed chromosomal target site between an $\alpha/\beta$-hydrolase gene fragment and an FMN-dependent NADH-azoreductase gene fragment. The integron carries \textit{aadB} ($=$\textit{ant(2$''$)-Ia}), \textit{aadA2}, and \textit{sul1}, plus the canonical class-1 3$'$-CS \textit{qacE$\Delta$1} and \textit{intI1}. AbGRI4 was absent from ABUH773.

\section{Claims and results}

\small
\begin{longtable}{@{}c p{7.2cm} c c@{}}
\toprule
\# & Claim & Tested? & Result \\
\midrule
\endhead
C1  & 11 replicons CP035043--CP035053 deposited and retrievable                           & \checkmark & Sizes match paper Table~1 exactly \\
C2  & All 4 isolates are Pasteur ST2                                                       & \checkmark & 4/4 ST2, allele profile 2-2-2-2-2-2-2 \\
C3  & All 4 isolates are Oxford ST281                                                      & partial   & Profile 1-17-\{3,189\}-2-2-99-3 (mlst DB lag) \\
C4  & AbGRI4 present at Table-1 coordinates in ABUH763/793/796                             & \checkmark & Three regions, each exactly 8{,}840 bp \\
C5  & AbGRI4 sequence essentially identical across the 3 positives                         & \checkmark & 0 mismatches / 8{,}840 bp all pairs \\
C6  & ABUH773 does NOT carry AbGRI4                                                        & \checkmark & Whole-genome AMR: 0 hits to aadB/aadA2/sul1/qacE$\Delta$1/intI1 \\
C7  & Integron carries aadB, aadA2, sul1                                                   & \checkmark & 3/3 databases at $\geq$99.87\% ID, $\geq$97.9\% length \\
C8  & Canonical class-1 3$'$-CS \textit{qacE$\Delta$1} + \textit{intI1}                    & \checkmark & \textit{qacE$\Delta$1} 100\% ID (CARD); \textit{intI1} in GBK \\
C9  & IS26 flanks at BOTH ends of the integron                                             & \checkmark & IS26 CDS at 5$'$ and 3$'$ in all 3 chromosomes \\
C10 & Insertion between EP550\_07220 ($\alpha/\beta$-hydrolase) and EP550\_07290 (azoreductase) & \checkmark & Both tags present, marked \texttt{/pseudo} in CP035043 \\
C11 & Target site is genuinely novel (not the canonical \textit{A.\ baumannii} target)      & \checkmark & Azo hits AB0057/17978 at $\geq$92.8\%; $\alpha/\beta$-hydrolase does NOT \\
C12 & Hybrid assembly required for IS-bounded RIs                                          & ---       & Out of scope (methods assertion) \\
\bottomrule
\end{longtable}
\normalsize

Testable claims: 11/12. Reproduced: 10/11 (+1 partial). No contradictions.

\section{Method}
All work on \texttt{uicgpu}, \texttt{/data/stevens/bvbrc69-abgri4/}, conda env \texttt{bvbrc14}.

\begin{enumerate}[leftmargin=*]
\item \textbf{Retrieval.} \texttt{efetch} (EDirect 24.0) --- 11 CP-numbered records deposited by the paper plus CP001182.2 (AB0057) and CP000521.1 (ATCC~17978). Script: \texttt{work/download\_genomes.sh}.
\item \textbf{Genome stats.} Biopython 1.87; per-replicon length and GC\%. Script: \texttt{work/genome\_stats.py}.
\item \textbf{MLST.} \texttt{mlst 2.33.1} (Torsten Seemann) with schemes \texttt{abaumannii\_2} (Pasteur) and \texttt{abaumannii} (Oxford). Script: \texttt{work/run\_mlst.sh}.
\item \textbf{AbGRI4 extraction.} Biopython slice at Table-1 coordinates $\pm$0. Script: \texttt{work/extract\_abgri4.py}.
\item \textbf{AMR annotation.} \texttt{abricate 1.4.0} against ResFinder, CARD, NCBI AMRFinderPlus, PlasmidFinder; \texttt{--minid 90 --mincov 80}. Script: \texttt{work/annotate\_abgri4.sh}.
\item \textbf{Identity / IS26 / target-site.} Pairwise Hamming (Biopython); GBK feature walk for IS26 and pseudogene flanks; direct locus\_tag lookup EP550\_07220/EP550\_07290 in CP035043.
\item \textbf{Comparative BLAST.} NCBI BLAST+ v2.16.0; \texttt{-evalue 1e-30 -perc\_identity 90} of flanks and full island against AB0057 and ATCC~17978; 20-kb window inspection around the AB0057 azoreductase hit.
\end{enumerate}

\section{Results (selected)}

\subsection{Genome statistics}
\begin{tabular}{@{}lrrrr@{}}\toprule
Strain & Chromosome & Length (bp) & GC\% & Total (bp) \\ \midrule
ABUH763 & CP035051 & 3{,}929{,}411 & 39.13 & 4{,}014{,}469 \\
ABUH773 & CP035049 & 3{,}873{,}900 & 39.06 & 3{,}885{,}710 \\
ABUH793 & CP035045 & 3{,}915{,}869 & 39.13 & 4{,}107{,}868 \\
ABUH796 & CP035043 & 3{,}930{,}797 & 39.12 & 3{,}943{,}749 \\
\bottomrule
\end{tabular}

\subsection{AbGRI4 span and identity}
\begin{tabular}{@{}lllr@{}}\toprule
Strain & Coordinates & Length & Hamming vs ABUH796 \\ \midrule
ABUH763 & CP035051:1518797--1527636       & 8{,}840 bp & 0/8840 \\
ABUH793 & CP035045:2219263--2228102 (rev) & 8{,}840 bp & 0/8840 (rev-comp) \\
ABUH796 & CP035043:1515737--1524576       & 8{,}840 bp & reference \\
\bottomrule
\end{tabular}

\subsection{AMR triad in AbGRI4 (ABUH796; ResFinder / CARD / NCBI, \%ID)}
\begin{tabular}{@{}lrccc@{}}\toprule
Gene & Position & ResFinder & CARD & NCBI \\ \midrule
\textit{ant(2$''$)-Ia (aadB)} & 2893--3426 (+) & 100.00 & 100.00 & 100.00 \\
\textit{aadA2}                 & 3474--4275 (+) & 99.88  & 99.87  & 100.00 \\
\textit{qacE$\Delta$1}         & 4439--4786 (+) & ---    & 100.00 & --- \\
\textit{sul1}                  & 4753--5619 (+) & 99.89  & 100.00 & 100.00 \\
\bottomrule
\end{tabular}

\subsection{Novel target site (BLAST vs external references)}
\begin{tabular}{@{}lll@{}}\toprule
Flank & AB0057 (CP001182.2) & ATCC 17978 (CP000521.1) \\ \midrule
EP550\_07290 (3$'$ azo, 309 bp)  & 92.88\% ID, 100\% len @ 1{,}646{,}616 & 93.20\% ID, 100\% len @ 1{,}586{,}111 \\
EP550\_07220 (5$'$ $\alpha/\beta$-h, 375 bp) & \textbf{no hit at $\geq$90\%}      & \textbf{no hit at $\geq$90\%} \\
\bottomrule
\end{tabular}

The 20-kb window around the AB0057 azoreductase (\texttt{AB57\_1564}) shows a LysR-family regulator (\texttt{AB57\_1563}) directly upstream --- \emph{not} an $\alpha/\beta$-hydrolase --- confirming the target-site gene pair is novel.

\section{Verdict: REPLICATED}
All 11 testable claims reproduced on real public data with standard bacterial-genomics tools; 3 independent AMR databases agree at $\geq$99.87\% ID; the negative-control isolate (ABUH773) is negative; the target site is genuinely novel with respect to the two canonical reference chromosomes. The one non-tested item (C12) is a methods assertion, not a data claim.

\section{GENUINE CRITIQUE}

This section records limitations that survived the replication --- items where the paper's story is consistent with what we can independently observe, but where a stronger claim would need work that was \emph{not} done here or is not decidable from the deposited artifacts alone.

\subsection*{1. Oxford ST281 is not confirmed by a modern MLST DB}
The current \texttt{mlst 2.33.1} \texttt{abaumannii} scheme returns a double-hit at \textit{gdhB} ($\{3, 189\}$) and consequently does not emit a single ST call. The allele profile 1-17-\{3,189\}-2-2-99-3 is \emph{consistent with} the paper's ST281 but not \emph{equal to} it as declared. A rigorous confirmation would require a fresh PubMLST query, manual allele resolution, or a version-pinned older scheme --- none of which were done. C3 is honestly marked ``partial''.

\subsection*{2. Byte-identity of the three islands is suspiciously perfect}
Zero mismatches over 8{,}840~bp across three independent clinical isolates is extraordinarily strong evidence of a single, recent horizontal-transfer event --- but it also lies at exactly the level of resolution where an assembly artifact would produce the same signal (e.g.\ if the same contig was propagated during hybrid polishing or if IS26-mediated collapse produced identical consensus). We did not re-assemble from raw reads; we accepted the deposited chromosomes. A stronger claim would independently re-polish from the SRA reads.

\subsection*{3. Class-1 integron reading is limited to the AMR-cassette triad}
We annotated \textit{aadB, aadA2, sul1, qacE$\Delta$1, intI1}. We did \emph{not} independently type the integron's $P_c$ promoter variant, \textit{attC} sites, or upstream Tn\textit{402} remnants; these features are relevant to expression strength and to placing AbGRI4 in the broader class-1 integron family tree but were beyond scope.

\subsection*{4. Novelty of the target site is bounded by our two comparators}
``Genuinely novel'' as demonstrated here means: not present at the canonical position in AB0057 or ATCC~17978. The paper's global-scale claim (AbGRI4 seen in non-ST2 lineages, suggesting HGT) was \emph{not} independently tested; we did not scan a broader panel of \textit{A.\ baumannii} genomes for the target-site gene pair or for AbGRI4 variants. A cross-population claim needs a cross-population dataset.

\subsection*{5. Absence in ABUH773 is a whole-genome AMR-panel result, not a coordinate check}
We ruled out AbGRI4 in ABUH773 by absence of the entire cassette on whole-genome ResFinder. A stricter test would examine the orthologous chromosomal region in ABUH773 for an unbroken $\alpha/\beta$-hydrolase/azoreductase pair (the ``empty'' target) --- confirming not just that the cassette is missing but that the target locus is intact. This is a small future check.

\subsection*{6. C12 (hybrid assembly requirement) is untested}
This is a methods observation. Independent literature supports it, but this replication does not.

\subsection*{7. Reproducibility of database versions}
Our AMR calls are pinned to abricate 1.4.0 + ResFinder (2026-Apr-3, 3206~seqs) / CARD (6052) / NCBI (8232) / PlasmidFinder (488). The paper used older databases in 2020. Differences in \%ID and allele names between our run and theirs are within the expected drift of database updates and do not affect the qualitative conclusion.

\section*{Provenance}
Scripts, raw \texttt{abricate} TSVs, GBK feature dumps, BLAST results, and extracted island FASTAs are archived under \texttt{report/evidence/} and \texttt{work/}.

\end{document}
